\chapter{Preface}

\section*{What these notes are}

Visual-inertial odometry (VIO) estimates the motion of a rigid body, such as a drone, a phone, a
headset or a car, by fusing an inertial measurement unit (IMU) with a camera. It is the
navigation core of most modern augmented-reality devices and of many autonomous aerial
vehicles. VIO is also one of the best ways to learn state estimation properly, because it forces
you to deal with almost every hard topic in the field at once: nonlinear dynamics, estimation on
curved spaces, sensor noise modeling, multi-rate fusion, unobservable directions, outliers,
numerical conditioning, and real-time computation.

The notes are written in the style of a graduate course. Each chapter states its learning
objectives, develops the theory with definitions, propositions and proofs, illustrates it with
worked examples and figures, flags the pitfalls that cost practitioners weeks of debugging, and
closes with key takeaways and exercises. All formulas and conventions match the \texttt{vio}
project specification in \texttt{AGENTS.md}. Every identity, Jacobian and growth law stated here has
been checked numerically, and the figures are generated by \texttt{scripts/make\_figures.py}, so you
can reproduce them and change their parameters.

\section*{Who should read them}

The assumed background is
\begin{itemize}
  \item multivariable calculus (Taylor series, chain rule, partial derivatives);
  \item introductory linear algebra (matrices, determinants, eigenvalues), which
        \cref{ch:linalg} revisits in the depth state estimation needs;
  \item basic probability (mean, variance, Gaussian distribution), revisited in \cref{ch:prob};
  \item for the implementation notes, working knowledge of modern C++ and Eigen.
\end{itemize}
No prior exposure to Lie groups, Kalman filtering or multi-view geometry is assumed.

\section*{How the notes are organized}

\begin{figure}[h]
\centering
\begin{tikzpicture}[node distance=5mm and 6mm, every node/.style={block, minimum width=27mm, font=\footnotesize}]
  \node (c1) {1 Introduction};
  \node (c2) [below left=6mm and 2mm of c1] {2 Linear algebra};
  \node (c3) [below right=6mm and 2mm of c1] {3 Probability};
  \node (c4) [below=6mm of c2] {4 Rotations};
  \node (c5) [below=6mm of c4] {5 Lie theory};
  \node (c6) [below=6mm of c3] {6 Kalman filter};
  \node (c7) [below=6mm of c6] {7 Error-state KF};
  \node (c8) [below left=8mm and -6mm of c5] {8 IMU};
  \node (c9) [right=4mm of c8] {9 Strapdown};
  \node (c10) [right=4mm of c9] {10 Error dynamics};
  \node (c11) [below=6mm of c10] {11 Updates};
  \node (c12) [right=4mm of c11] {12 Consistency};
  \node (c13) [below=6mm of c8] {13 Camera};
  \node (c14) [below=6mm of c13] {14 Multi-view};
  \node (c15) [below=6mm of c14] {15 Front end};
  \node (c16) [below=6mm of c11] {16 MSCKF idea};
  \node (c17) [below=6mm of c16] {17 Cloning};
  \node (c18) [below=6mm of c17] {18 MSCKF update};
  \node (c19) [right=4mm of c18] {19 Full algorithm};
  \node (c20) [above=6mm of c19, fill=white] {20 Observability};
  \node (c21) [below=6mm of c19, fill=white] {21 Engineering};
  \draw[->, primary] (c1) -- (c2); \draw[->, primary] (c1) -- (c3);
  \draw[->, primary] (c2) -- (c4); \draw[->, primary] (c4) -- (c5);
  \draw[->, primary] (c3) -- (c6); \draw[->, primary] (c6) -- (c7);
  \draw[->, primary] (c5) -- (c7);
  \draw[->, primary] (c5) -- (c8); \draw[->, primary] (c8) -- (c9); \draw[->, primary] (c9) -- (c10);
  \draw[->, primary] (c7) -- (c10); \draw[->, primary] (c10) -- (c11); \draw[->, primary] (c11) -- (c12);
  \draw[->, primary] (c8) -- (c13); \draw[->, primary] (c13) -- (c14); \draw[->, primary] (c14) -- (c15);
  \draw[->, primary] (c11) -- (c16); \draw[->, primary] (c16) -- (c17); \draw[->, primary] (c17) -- (c18);
  \draw[->, primary] (c14) -- (c18); \draw[->, primary] (c18) -- (c19);
  \draw[->, primary] (c15) to[out=-10, in=200] (c19);
  \draw[->, primary] (c12) -- (c20); \draw[->, primary] (c19) -- (c21);
\end{tikzpicture}
\caption{Chapter dependencies. Arrows point from prerequisites to the chapters that use them.
Chapters 8--12 can be read before Part V, but the MSCKF chapters need both branches.}
\label{fig:dependencies}
\end{figure}

\begin{description}[leftmargin=0pt, itemsep=0.4em]
  \item[Part I, Mathematical Foundations.] The vocabulary of the whole subject: linear algebra and
    the factorizations that make filters numerically sound, probability and Gaussian algebra,
    rotations and quaternions, and Lie theory, the calculus of the rotation group.
  \item[Part II, Estimation Theory.] The Kalman filter from Bayes' rule, its nonlinear extension,
    and the error-state formulation that makes estimation on a manifold clean.
  \item[Part III, Inertial Navigation.] IMU physics and noise, strapdown propagation, and the
    linearized error dynamics that drive covariance growth.
  \item[Part IV, The Error-State Filter in Practice.] Measurement updates, gating, injection and
    reset, and the statistics that tell you whether the filter is lying to you.
  \item[Part V, Geometric Computer Vision.] Camera models, epipolar geometry, triangulation,
    corner detection and optical-flow tracking.
  \item[Part VI, The MSCKF.] Stochastic cloning, the visual measurement model, left-nullspace
    projection, and the complete algorithm.
  \item[Part VII, Beyond the Baseline.] Observability and consistency, and engineering a
    deterministic real-time implementation in C++20.
\end{description}

\section*{Reading the boxes}

Five kinds of callout recur throughout:

\begin{keyidea}[the one-sentence version]
The idea to remember even if you forget the derivation.
\end{keyidea}
\begin{convention}[project specification]
A rule from \texttt{AGENTS.md} that is frozen for the whole project. These are never changed
silently. A change is a logged design decision.
\end{convention}
\begin{pitfall}[what goes wrong in practice]
A mistake that is easy to make, hard to see, and expensive to debug.
\end{pitfall}
\begin{implnote}[from math to code]
How an equation becomes deterministic, allocation-free C++20 with Eigen.
\end{implnote}
\begin{projectlink}[Milestone N]
Where the material meets a project milestone and its exit criteria.
\end{projectlink}

A sixth box, the \textbf{Derivation exercise}, marks the Jacobians that the project asks you to
derive yourself: $F_c$, $G_c$, the measurement Jacobians $H$, and the cloning Jacobian $J$. The
surrounding text teaches the method on closely related examples, and the box gives hints plus a
finite-difference test that tells you when your answer is right. The worked solutions are
collected in an appendix that exists only in the \emph{solutions} edition of these notes\ifsolutions{} (this one: \cref{app:solutions})\else{} (this edition omits it)\fi.

\section*{Milestones and chapters}

\begin{table}[h]
\centering\small
\begin{tabularx}{\textwidth}{@{}l X l@{}}
\toprule
\textbf{Milestone} & \textbf{Topic} & \textbf{Chapters} \\
\midrule
M1 & Build harness, sanitizers, allocation guard & \ref{ch:engineering} \\
M2 & Lie group math on $\SO$ & \ref{ch:rotations}, \ref{ch:lie} \\
M3 & EuRoC dataset streamer & \ref{ch:imu}, \cref{app:euroc} \\
M4 & Nominal IMU kinematics & \ref{ch:strapdown} \\
M5 & Error covariance propagation & \ref{ch:errordyn} \\
M6 & Pose update, injection and reset & \ref{ch:eskf}, \ref{ch:updates} \\
M7 & Telemetry and consistency & \ref{ch:consistency} \\
M8 & Feature detection and optical flow & \ref{ch:camera}, \ref{ch:frontend} \\
M9 & Epipolar geometry and outlier rejection & \ref{ch:multiview}, \ref{ch:frontend} \\
M10a & Stochastic cloning and sliding window & \ref{ch:msckfintro}, \ref{ch:cloning} \\
M10b & MSCKF visual update & \ref{ch:msckfupdate}, \ref{ch:fullalgo}, \ref{ch:observability} \\
\bottomrule
\end{tabularx}
\end{table}

\section*{Sources and acknowledgments}

These notes stand on a body of excellent literature. The treatment of quaternions and the
error-state filter follows the conventions of Solà~\cite{sola2017quaternion}; Lie theory follows
Solà, Deray and Atchuthan~\cite{sola2018micro} and Barfoot~\cite{barfoot2017state}; the MSCKF is
due to Mourikis and Roumeliotis~\cite{mourikis2007msckf} and was refined by Li and
Mourikis~\cite{li2013msckf2}; much practical wisdom comes from OpenVINS~\cite{geneva2020openvins}
and MSCKF-VIO~\cite{sun2018msckfvio}. The annotated bibliography at the end points to the best
places to go deeper on each topic.
